% =============================================================================
%  Weeks 3-4 Progress Report
%  Probabilistic Motor Subtype Progression Model for Parkinson's Disease
%
%  Compiles with pdfLaTeX. Figures are read from ../outputs/transitions/figures,
%  regenerated by  python scripts/02_transition_analysis.py
% =============================================================================
\documentclass[11pt,a4paper]{article}

\usepackage[T1]{fontenc}
\usepackage[utf8]{inputenc}
\usepackage[margin=2.4cm]{geometry}
\usepackage{booktabs}
\usepackage{longtable}
\usepackage{array}
\usepackage{xcolor}
\usepackage{graphicx}
\usepackage{enumitem}
\usepackage{caption}
\usepackage{amsmath}
\usepackage{amssymb}
\usepackage[hidelinks]{hyperref}

\graphicspath{{../outputs/transitions/figures/}}

\definecolor{accent}{HTML}{1F4E4C}
\definecolor{warn}{HTML}{B7770B}
\definecolor{keybg}{HTML}{F4F6F8}

\newcommand{\code}[1]{\texttt{\small #1}}
\newcommand{\arr}{$\rightarrow$}
\newcommand{\TD}{TD}
\newcommand{\IND}{Indeterminate}
\newcommand{\PIGD}{PIGD}

\setlist[itemize]{itemsep=2pt,topsep=4pt}
\setlist[enumerate]{itemsep=2pt,topsep=4pt}
\captionsetup{font=small,labelfont=bf}

% A ruled "key point" block for the things worth memorising. Rules rather than
% a colorbox, because a colorbox cannot span an environment boundary.
\newenvironment{keypoint}
  {\begin{center}\begin{minipage}{0.94\textwidth}%
   \vspace{4pt}\hrule\vspace{7pt}\small}
  {\vspace{7pt}\hrule\vspace{4pt}\end{minipage}\end{center}}

\title{\textbf{Longitudinal Analysis of Motor Subtype Transitions\\ in Parkinson's Disease}\\[0.6em]
       \large Weeks 3--4 Progress Report}
\author{Parv\\[0.3em]
        \small Supervisor: Prof.\ K.~R.\ Sahasranand\\
        \small Department of Data Science, IIT Palakkad}
\date{August 2026}

\begin{document}
\maketitle

\begin{abstract}
\noindent
This report covers Weeks 3--4 of the project: tracking TD/PIGD motor subtype
changes across visits, estimating transition probabilities over different
follow-up intervals, and visualising progression pathways. The analysis uses
3{,}361 consecutive labelled visit pairs from 934 sporadic PD patients in the
PPMI cohort.

Four results shape the modelling that follows. \emph{Indeterminate} is not a
stable state but a crossing --- only 19.2\% of patients remain in it at the
next visit. \emph{PIGD is not absorbing}: 29.8\% of PIGD visits are followed by
a non-PIGD assessment. \emph{35.2\% of all observed subtype changes reverse at
the very next visit}, so a transition matrix built from raw labels measures
assessment noise as much as disease progression. Yet the movement is not
symmetric: TD~\arr~PIGD outnumbers PIGD~\arr~TD by 253 to 141
($p = 1.8\times10^{-8}$), so a genuine directional drift exists beneath the
noise. Finally, transition probabilities scale strongly with the gap between
visits, which means a single pooled transition matrix is not a well-defined
quantity and the Week 5--7 model should work in continuous time.
\end{abstract}

\vspace{0.4em}

% =============================================================================
\section{Objective and scope}
% =============================================================================

The project proposal assigns Weeks 3--4 the following tasks:

\begin{itemize}
  \item Track subtype changes for each patient across visits.
  \item Compute transition frequencies (TD~\arr~TD, TD~\arr~Indeterminate,
        TD~\arr~PIGD, and so on).
  \item Estimate transition probabilities over different follow-up intervals.
  \item Visualise common progression pathways.
\end{itemize}

\noindent
with the deliverable stated as \emph{transition matrices and visualisations of
subtype progression}. All four are complete. This report also records two
analyses that were not requested but that materially change how the transition
matrix should be read --- a reversion analysis and a sustained-conversion
analysis --- because without them the headline numbers are misleading.

Everything here is descriptive. No model is fitted; that is Weeks 5--7.

% =============================================================================
\section{Data}
% =============================================================================

\subsection{The panel}

The input is \code{outputs/visit\_panel.csv}, built in Weeks 1--2 from the PPMI
Curated Data Cut joined to the item-level MDS-UPDRS Part II and Part III files.
Subtype labels follow Stebbins et al.\ (2013):

\[
  \text{ratio} = \frac{\text{mean of 11 tremor items}}{\text{mean of 5 gait/postural items}}
\]

with $\text{ratio} \geq 1.15 \Rightarrow$ TD, $\text{ratio} \leq 0.90
\Rightarrow$ PIGD, and the interval between them Indeterminate. Scores are taken
from the OFF-medication or untreated examination.

\begin{table}[htbp]
\centering
\begin{tabular}{@{}lr@{}}
\toprule
\textbf{Quantity} & \textbf{Value} \\
\midrule
Patients (sporadic PD)                       & 1{,}080 \\
Visits                                       & 5{,}774 \\
Visits with a subtype label                  & 4{,}441 \quad (76.9\%) \\
Patients contributing $\geq 1$ transition    & 934 \\
Consecutive labelled visit pairs             & 3{,}361 \\
Median interval between labelled visits      & 1.00 years \\
Maximum follow-up                            & 14 years \\
\bottomrule
\end{tabular}
\caption{Analysis population.}
\end{table}

\subsection{Why the labels can be trusted}

The recomputed labels were validated against both of PPMI's own derivations
before any of this analysis was run:

\begin{table}[htbp]
\centering
\begin{tabular}{@{}llr@{}}
\toprule
\textbf{Reference} & \textbf{Classes} & \textbf{Agreement} \\
\midrule
Curated cut \code{td\_pigd}      & 2 (TD / non-TD)             & 4{,}431 / 4{,}431 \; = \; 100.00\% \\
Archived \code{td\_pigd\_old}    & 3 (TD / Indet.\ / PIGD)     & 1{,}595 / 1{,}597 \; = \; 99.87\% \\
\bottomrule
\end{tabular}
\caption{Label validation. The three-class check covers 390 patients and is the
only available test of the Indeterminate/PIGD boundary, since the curated cut
collapses those two classes.}
\end{table}

\subsection{Baseline distribution}

At baseline: TD 71.7\%, Indeterminate 10.1\%, PIGD 18.1\%. Pooled over all
labelled visits: TD 65.1\%, Indeterminate 9.9\%, PIGD 25.0\%.

% =============================================================================
\section{Methods}
% =============================================================================

\subsection{Transition pairs}

For each patient, visits with a non-missing label are sorted by elapsed time
and adjacent pairs are formed. Two decisions matter:

\begin{enumerate}
  \item \textbf{Unlabelled visits are skipped, not treated as terminal.} If a
        patient is labelled at years 0 and 2 but not at year 1, that yields one
        pair with a two-year interval rather than truncating the sequence.
  \item \textbf{The elapsed interval is carried through.} PPMI visits are
        annual early on and biennial later, and patients miss visits. A pair is
        therefore not a fixed unit of time.
\end{enumerate}

Time is measured from the patient's own first visit using \code{INFODT}, which
PPMI releases at month resolution.

\subsection{Transition matrix and uncertainty}

The matrix is row-stochastic: entry $(i,j)$ is the proportion of visits in
state $i$ followed by state $j$. Row counts differ by an order of magnitude
(2{,}261 from TD against 334 from Indeterminate), so every probability is
reported with a \textbf{Wilson score interval} rather than a point estimate
alone. Wilson is used in preference to the normal approximation because the
latter produces intervals that leave $[0,1]$ for probabilities near the
boundary, which several of these are.

\subsection{Interval stratification}

Pairs are binned by elapsed interval ($\leq 1.5$, $1.5$--$2.5$, $2.5$--$4.0$,
$> 4.0$ years) and a separate matrix estimated in each bin.

\subsection{Reversion analysis}

For every observed change of state $A \rightarrow B$ where a third labelled
visit exists, we record whether that third visit returns to $A$. This separates
two situations a transition matrix reports identically: a patient whose
phenotype genuinely progressed, and a patient whose single noisy assessment
crossed a threshold and crossed back.

\subsection{Time to conversion}

Time from baseline to a patient's first PIGD assessment is estimated by
Kaplan--Meier with Greenwood standard errors, censoring at the last labelled
visit. Patients already PIGD at baseline are excluded from the risk set. Two
definitions of conversion are compared:

\begin{itemize}
  \item \textbf{naive} --- a single PIGD assessment;
  \item \textbf{sustained} --- PIGD at two consecutive labelled visits.
\end{itemize}

\subsection{Pathways}

Each patient's label sequence is reduced by collapsing immediate repeats, so
(TD, TD, PIGD, PIGD) becomes TD~\arr~PIGD. Without this,
pathway frequencies would be dominated by how often a patient attended rather
than by what happened to them.

% =============================================================================
\section{Results}
% =============================================================================

\subsection{Pooled transition matrix}

\begin{table}[htbp]
\centering
\begin{tabular}{@{}l rrr r@{}}
\toprule
 & \multicolumn{3}{c}{\textbf{State at next visit}} & \\
\cmidrule(lr){2-4}
\textbf{Current state} & \textbf{TD} & \textbf{Indeterminate} & \textbf{PIGD} & \textbf{Row $n$} \\
\midrule
TD             & \textbf{0.808} & 0.080 & 0.112 & 2{,}261 \\
Indeterminate  & 0.440 & \textbf{0.192} & 0.368 & 334 \\
PIGD           & 0.184 & 0.114 & \textbf{0.702} & 766 \\
\bottomrule
\end{tabular}
\caption{Transition probabilities between consecutive labelled visits
($n = 3{,}361$ pairs). Counts: TD row 1827/181/253; Indeterminate row
147/64/123; PIGD row 141/87/538.}
\label{tab:matrix}
\end{table}

\begin{table}[htbp]
\centering
\small
\begin{tabular}{@{}ll rr r l@{}}
\toprule
\textbf{From} & \textbf{To} & \textbf{$n$ from} & \textbf{$n$ trans.} & \textbf{Prob.} & \textbf{Wilson 95\% CI} \\
\midrule
TD            & TD            & 2261 & 1827 & 0.808 & [0.791, 0.824] \\
TD            & Indeterminate & 2261 &  181 & 0.080 & [0.070, 0.092] \\
TD            & PIGD          & 2261 &  253 & 0.112 & [0.100, 0.126] \\
Indeterminate & TD            &  334 &  147 & 0.440 & [0.388, 0.494] \\
Indeterminate & Indeterminate &  334 &   64 & 0.192 & [0.153, 0.237] \\
Indeterminate & PIGD          &  334 &  123 & 0.368 & [0.318, 0.421] \\
PIGD          & TD            &  766 &  141 & 0.184 & [0.158, 0.213] \\
PIGD          & Indeterminate &  766 &   87 & 0.114 & [0.093, 0.138] \\
PIGD          & PIGD          &  766 &  538 & 0.702 & [0.669, 0.734] \\
\bottomrule
\end{tabular}
\caption{Transition probabilities with Wilson score intervals.}
\end{table}

\begin{figure}[htbp]
\centering
\includegraphics[width=0.78\textwidth]{transition_matrix.png}
\caption{Pooled transition matrix. Colour encodes probability magnitude on a
single-hue sequential ramp; cell counts are printed because the row totals
differ by an order of magnitude.}
\end{figure}

\clearpage

\subsection{Persistence of each state}

From the diagonal, the expected number of consecutive visits spent in a state
before leaving it is $1/(1-p_{\text{stay}})$:

\begin{table}[htbp]
\centering
\begin{tabular}{@{}lrr@{}}
\toprule
\textbf{State} & \textbf{$p_{\text{stay}}$} & \textbf{Mean consecutive visits} \\
\midrule
TD             & 0.808 & 5.21 \\
PIGD           & 0.702 & 3.36 \\
Indeterminate  & 0.192 & 1.24 \\
\bottomrule
\end{tabular}
\caption{State persistence. At a median interval of one year these are roughly
years. Indeterminate barely outlives a single visit.}
\end{table}

\subsection{Directionality: is the movement symmetric?}

High reversion could mean the labels are pure noise. Testing each pair of
states against the null of symmetric exchange (exact binomial):

\begin{table}[htbp]
\centering
\begin{tabular}{@{}lrrrl@{}}
\toprule
\textbf{Pair} & \textbf{Forward} & \textbf{Reverse} & \textbf{Net} & \textbf{$p$} \\
\midrule
TD $\rightarrow$ PIGD          & 253 & 141 & $+112$ & $1.8\times10^{-8}$ \\
Indeterminate $\rightarrow$ PIGD & 123 &  87 & $+36$  & $0.016$ \\
TD $\rightarrow$ Indeterminate & 181 & 147 & $+34$  & $0.068$ \\
\bottomrule
\end{tabular}
\caption{Directional asymmetry. Of 932 cross-state moves, only 182 --- 19.5\%
--- are net drift; the remaining 80.5\% cancels out. But the residue is real
and, for TD $\leftrightarrow$ PIGD, overwhelmingly significant.}
\label{tab:symmetry}
\end{table}

\begin{keypoint}
\textbf{The single most important sentence in this report.} Most observed
subtype movement is bidirectional churn, but it does not cancel exactly: there
is a statistically overwhelming net drift from TD toward PIGD. The data is
therefore neither ``clean progression'' nor ``pure noise'' --- it is a genuine
drift observed through a noisy measurement. That is precisely the structure a
latent-state model with a measurement-error term is built to represent, and it
is the empirical justification for the proposed approach.
\end{keypoint}

\subsection{Dependence on the follow-up interval}

\begin{table}[htbp]
\centering
\small
\begin{tabular}{@{}l rrr rrr rrr@{}}
\toprule
& \multicolumn{3}{c}{\textbf{From TD}} & \multicolumn{3}{c}{\textbf{From Indeterminate}} & \multicolumn{3}{c}{\textbf{From PIGD}} \\
\cmidrule(lr){2-4}\cmidrule(lr){5-7}\cmidrule(lr){8-10}
\textbf{Interval (yr)} & TD & Ind & PIGD & TD & Ind & PIGD & TD & Ind & PIGD \\
\midrule
$\leq 1.5$   & 0.820 & 0.079 & 0.100 & 0.477 & 0.197 & 0.326 & 0.192 & 0.115 & 0.693 \\
$1.5$--$2.5$ & 0.762 & 0.069 & 0.168 & 0.200 & 0.200 & 0.600 & 0.100 & 0.100 & 0.800 \\
$2.5$--$4.0$ & 0.750 & 0.097 & 0.153 & 0.375 & 0.125 & 0.500 & 0.174 & 0.130 & 0.696 \\
$> 4.0$      & 0.578 & 0.133 & 0.289 & 0.222 & 0.111 & 0.667 & 0.250 & 0.083 & 0.667 \\
\bottomrule
\end{tabular}
\caption{Transition probabilities stratified by elapsed interval. Row $n$, in
the same interval order, is 1942/202/72/45 from TD, 279/30/16/9 from
Indeterminate and 661/70/23/12 from PIGD. The Indeterminate and long-interval
cells rest on very few observations and should be read as intervals, not point
estimates.}
\end{table}

\begin{figure}[htbp]
\centering
\includegraphics[width=0.88\textwidth]{transitions_by_interval.png}
\caption{Selected transitions against elapsed interval. TD~\arr~PIGD nearly
triples from 0.100 to 0.289 as the gap widens. Washed-out bars marked $\dagger$
rest on fewer than 30 transitions.}
\end{figure}

\clearpage

\subsection{Reversion: how much of this is noise?}

\begin{table}[htbp]
\centering
\begin{tabular}{@{}ll rr r@{}}
\toprule
\textbf{From} & \textbf{To} & \textbf{$n$ changes} & \textbf{$n$ reverted} & \textbf{Reversion rate} \\
\midrule
TD            & Indeterminate & 134 & 75 & \textbf{0.560} \\
PIGD          & Indeterminate &  54 & 27 & \textbf{0.500} \\
PIGD          & TD            & 112 & 44 & 0.393 \\
TD            & PIGD          & 172 & 55 & 0.320 \\
Indeterminate & PIGD          &  93 & 17 & 0.183 \\
Indeterminate & TD            & 108 & 19 & 0.176 \\
\midrule
\multicolumn{2}{@{}l}{\textbf{Overall}} & \textbf{673} & \textbf{237} & \textbf{0.352} \\
\bottomrule
\end{tabular}
\caption{Share of observed state changes that return to the origin state at the
very next visit.}
\label{tab:reversion}
\end{table}

\begin{figure}[htbp]
\centering
\includegraphics[width=0.80\textwidth]{reversion.png}
\caption{Reversion rates. Changes \emph{into} Indeterminate reverse most often
(50--56\%), consistent with Indeterminate being a threshold band that a noisy
ratio wanders across rather than a clinical state a patient occupies.}
\end{figure}

\subsection{Time to PIGD conversion}

\begin{table}[htbp]
\centering
\begin{tabular}{@{}lrr@{}}
\toprule
 & \textbf{Naive (1 visit)} & \textbf{Sustained (2 visits)} \\
\midrule
At risk (not PIGD at baseline) & 885 & 885 \\
Converted                      & 263 \; (29.7\%) & 127 \; (14.4\%) \\
Median years to conversion     & 3.00 & 3.00 \\
KM cumulative incidence, year 2 & 12.5\% & 6.5\% \\
KM cumulative incidence, year 5 & 35.3\% & 15.3\% \\
\bottomrule
\end{tabular}
\caption{Requiring the conversion to hold for two consecutive visits removes
136 of 263 apparent conversions --- 51.7\%.}
\label{tab:conversion}
\end{table}

\begin{figure}[htbp]
\centering
\includegraphics[width=0.86\textwidth]{time_to_pigd_conversion.png}
\caption{Cumulative incidence of PIGD conversion under the two definitions.
Curves are truncated at year 12, where the risk set falls below 25; the
at-risk counts are printed beneath the axis.}
\end{figure}

\clearpage

\subsection{Prevalence across follow-up}

\begin{table}[htbp]
\centering
\small
\begin{tabular}{@{}r rrr rrr r@{}}
\toprule
\textbf{Year} & \textbf{$n$ TD} & \textbf{$n$ Ind} & \textbf{$n$ PIGD} & \textbf{\% TD} & \textbf{\% Ind} & \textbf{\% PIGD} & \textbf{$n$ total} \\
\midrule
0  & 772 & 109 & 195 & 71.7 & 10.1 & 18.1 & 1076 \\
1  & 556 &  84 & 166 & 69.0 & 10.4 & 20.6 &  806 \\
2  & 467 &  59 & 151 & 69.0 &  8.7 & 22.3 &  677 \\
3  & 332 &  39 & 109 & 69.2 &  8.1 & 22.7 &  480 \\
4  & 209 &  30 &  90 & 63.5 &  9.1 & 27.4 &  329 \\
5  & 129 &  27 &  60 & 59.7 & 12.5 & 27.8 &  216 \\
6  & 102 &  16 &  64 & 56.0 &  8.8 & 35.2 &  182 \\
7  &  94 &  19 &  64 & 53.1 & 10.7 & 36.2 &  177 \\
8  &  58 &  19 &  44 & 47.9 & 15.7 & 36.4 &  121 \\
9  &  50 &  11 &  38 & 50.5 & 11.1 & 38.4 &   99 \\
10 &  47 &  12 &  49 & 43.5 & 11.1 & 45.4 &  108 \\
11 &  36 &  11 &  45 & 39.1 & 12.0 & 48.9 &   92 \\
\bottomrule
\end{tabular}
\caption{Cross-sectional prevalence. Years 12--14 are omitted here
($n \leq 52$) but appear in \code{prevalence\_by\_year.csv}.}
\end{table}

\begin{figure}[htbp]
\centering
\includegraphics[width=0.86\textwidth]{prevalence_over_time.png}
\caption{Subtype prevalence with the assessed-$n$ panel beneath. The $n$ panel
is part of the reading: by year 8 fewer than 12\% of the baseline cohort is
still contributing.}
\end{figure}

\clearpage

\subsection{Progression pathways}

\begin{table}[htbp]
\centering
\begin{tabular}{@{}l rr@{}}
\toprule
\textbf{Pathway (repeats collapsed)} & \textbf{Patients} & \textbf{Share} \\
\midrule
TD                                       & 399 & 42.7\% \\
TD \arr\ PIGD                            &  83 &  8.9\% \\
PIGD                                     &  71 &  7.6\% \\
TD \arr\ Indeterminate                   &  39 &  4.2\% \\
TD \arr\ Indeterminate \arr\ TD          &  35 &  3.7\% \\
Indeterminate \arr\ PIGD                 &  27 &  2.9\% \\
Indeterminate \arr\ TD                   &  26 &  2.8\% \\
PIGD \arr\ TD                            &  24 &  2.6\% \\
TD \arr\ Indeterminate \arr\ PIGD        &  21 &  2.2\% \\
TD \arr\ PIGD \arr\ TD                   &  20 &  2.1\% \\
PIGD \arr\ TD \arr\ PIGD                 &  19 &  2.0\% \\
Indeterminate \arr\ TD \arr\ PIGD        &  11 &  1.2\% \\
\bottomrule
\end{tabular}
\caption{The twelve commonest of 86 distinct pathways across the 934 patients
with two or more labelled visits.}
\end{table}

\noindent
476 of those 934 patients (51.0\%) show no recorded change of state. Median
number of state changes per patient is 0; mean 1.00; maximum 10. A pathway
requires at least two labelled visits, so the 146 patients seen once are
excluded: counting them would credit them with a stable course they never had
the opportunity to depart from. Note also that \emph{no recorded change} is not
the same as \emph{no progression} --- a patient can worsen substantially without
ever crossing a subtype threshold.

\begin{figure}[htbp]
\centering
\includegraphics[width=0.94\textwidth]{pathway_flow.png}
\caption{Subtype pathways for the 274 patients assessed at baseline, year 2 and
year 4. Restricted to complete cases, so free of attrition but not of the
selection that surviving with a full record implies.}
\end{figure}

\clearpage

% =============================================================================
\section{What this means for the model}
% =============================================================================

\subsection{Four consequences}

\begin{enumerate}
  \item \textbf{Indeterminate should not be modelled as a clinical state.} With
        $p_{\text{stay}} = 0.192$ and reversion rates of 50--56\% for changes
        into it, it behaves exactly as the neutral band of a continuous latent
        variable would: a region the ratio passes through, not a place a
        patient sits. This is direct empirical support for the proposal's
        design, which models the continuous ratio and derives the three classes
        by thresholding, rather than modelling three discrete states.

  \item \textbf{The observation model needs a real noise term.} 35.2\%
        immediate reversion cannot be explained by disease progression. The
        proposal's $\text{observed} = \text{latent} + \text{noise}$ structure is
        not a modelling convenience here; it is required to fit what the data
        does.

  \item \textbf{The model must work in continuous time.} TD~\arr~PIGD nearly
        triples between the shortest and longest interval bins. A discrete
        Markov chain over ``consecutive visits'' would absorb PPMI's attendance
        schedule into its progression estimate. The natural fix is a
        continuous-time formulation --- estimate a generator matrix $Q$ and use
        $P(t) = e^{Qt}$ --- or, following the proposal, model the latent ratio
        as a function of actual elapsed time.

  \item \textbf{Conversion must be defined before it is predicted.} The
        naive/sustained gap (35.3\% versus 15.3\% at five years) is larger than
        any effect the model is likely to estimate. This choice should be made
        and stated now, not discovered later.
\end{enumerate}

\subsection{A caution on the pooled matrix}

Taking Table~\ref{tab:matrix} at face value and iterating it gives a stationary
distribution of TD 0.560, Indeterminate 0.103, PIGD 0.338. The observed year-11
prevalence is TD 0.391, Indeterminate 0.120, PIGD 0.489 --- considerably more
PIGD than the pooled matrix predicts. At least three things could explain the
gap: transition rates are not constant over disease duration, the process is
not Markov in the observed label, or attrition is selective. A fourth
possibility is developed in Section~\ref{sec:missingness} --- that the
transition matrix itself is estimated on a sample from which the most severely
affected visits have been systematically dropped, so its rates understate the
true drift. Disentangling these is a Week 5--7 question, but the discrepancy is
worth stating rather than presenting a stationary distribution as a finding.

% =============================================================================
\section{Limitations}
% =============================================================================

\subsection{Label availability is not random, and it is not neutral}
\label{sec:missingness}

Of the 5{,}774 cohort visits, 1{,}333 (23.1\%) carry no subtype label. That
missingness is not a random thinning of the data, and its direction matters.

\begin{table}[htbp]
\centering
\begin{tabular}{@{}lrl@{}}
\toprule
\textbf{Reason a visit has no label} & \textbf{$n$} & \textbf{Share of the loss} \\
\midrule
No OFF-medication Part III exam recorded  & 722 & 54.2\% \\
OFF exam recorded but items incomplete    & 595 & 44.6\% \\
Part II missing or incomplete             &  16 &  1.2\% \\
\bottomrule
\end{tabular}
\caption{Why visits go unlabelled.}
\end{table}

\noindent
Label availability also falls sharply across follow-up: 99.6\% of baseline
visits are labelled, 75.3\% at year 4, 55.4\% at year 11, and 19.2\% at
year 13.

\paragraph{The second route is the concerning one.} Among the 1{,}253
OFF-medication exams that were recorded but could not be scored,
\textbf{NP3PSTBL --- item 3.12, postural stability --- is missing in 1{,}243 of
them}, and 185 of those are explicitly coded \code{101}, ``Unable to Rate''. The
pull test is frequently unsafe or impossible to administer in a patient with
poor balance.

Because our rule requires all sixteen items, one unscored postural-stability
test discards the whole visit. And the visits it discards are not a random
sample: exams with an incomplete item set come from materially more advanced
patients, with a median Hoehn \& Yahr stage of 3.0 against 2.0 for complete
exams (Mann--Whitney $p = 1.4 \times 10^{-17}$).

\begin{keypoint}
\textbf{This biases the results in a knowable direction.} Postural stability is
one of the five \emph{PIGD} items. The visits we lose are therefore
disproportionately those of patients whose gait and balance are worst --- the
patients most likely to be classified PIGD. The observed drift toward PIGD, and
the PIGD conversion rates in Table~\ref{tab:conversion}, are for that reason
most likely \emph{underestimates}. This also offers a plausible explanation for
the gap noted in Section~5.2, where observed year-11 PIGD prevalence exceeds
what the pooled matrix predicts.
\end{keypoint}

\noindent
Two things are worth noting about how this was established. First, an earlier
reading of the same numbers attributed the loss mainly to ON-state exams
replacing OFF-state ones as patients become medicated. That explanation
accounts for only about half the loss and plateaus after year 5, whereas
incomplete item sets persist and are the clearly severity-linked route; the
severity comparison above is what distinguished them. Second, the effect is a
consequence of the strict all-items rule inherited from PPMI's own derivation.
Relaxing it to tolerate one missing item was tested in Weeks 1--2 and recovers
only 44 visits while changing no existing label, so relaxation is not a
remedy --- the incomplete exams are typically missing many items, not one.

A sensitivity analysis that imputes the postural-stability item, or that models
the label as missing-not-at-random, would quantify the size of this bias. That
is a Weeks 5--7 task.

\subsection{Other limitations}

\begin{itemize}
  \item \textbf{Sparse cells.} Transitions out of Indeterminate at intervals
        beyond 1.5 years rest on 9--30 observations. The Wilson intervals are
        reported for exactly this reason and those cells are flagged in the
        figure.

  \item \textbf{Selective attrition.} By year 8 fewer than 12\% of the baseline
        cohort is still assessed. This is separate from, and compounds, the
        label-availability problem in Section~\ref{sec:missingness}: a visit
        must both happen and be scorable to contribute.

  \item \textbf{Medication state.} Labels use the OFF or untreated exam, which
        is correct, but the proportion of patients who are genuinely untreated
        falls over follow-up --- from 99.7\% at baseline to under 2\% by year 7.
        Some of the observed churn may reflect changing treatment context
        rather than phenotype.

  \item \textbf{Month-resolution dates.} PPMI releases \code{INFODT} as month
        and year only, so intervals carry roughly $\pm$ two weeks of noise.
        Negligible at annual spacing.

  \item \textbf{The reversion analysis needs three visits.} Patients with
        exactly two labelled visits contribute to the transition matrix but not
        to Table~\ref{tab:reversion}, so the reversion estimate is drawn from
        the better-attended subset.

  \item \textbf{Ratio undefined for 16.3\% of labelled visits} (725 of 4{,}441)
        where the PIGD subscore is exactly zero. Those visits are classifiable
        --- the rule sends them to TD --- but carry no continuous observation.
        This is carried over from Weeks 1--2 and remains open.

  \item \textbf{Sporadic PD only.} Genetic PD subgroups (LRRK2, GBA, and
        others) are excluded. Extending to them is a configuration change.
\end{itemize}

% =============================================================================
\section{Anticipated questions}
% =============================================================================

\noindent\textbf{Q. If 35\% of changes reverse immediately, are the labels just noise?}\\
No. If the movement were pure noise it would be symmetric, and it is not:
TD~\arr~PIGD outnumbers PIGD~\arr~TD 253 to 141, $p = 1.8\times10^{-8}$
(Table~\ref{tab:symmetry}). Of 932 cross-state moves, 80.5\% cancels and 19.5\%
is net drift toward PIGD. So there is a real signal, observed through a noisy
measurement --- which is the case for a latent-variable model, not against it.

\medskip
\noindent\textbf{Q. Why is Indeterminate so unstable? Is that a bug?}\\
It is expected, and it was checked. Indeterminate is defined as
$0.90 < \text{ratio} < 1.15$ --- a narrow band on a continuous, noisily measured
quantity. Anyone whose true ratio sits near the band will cross in and out
between visits. The label logic was validated at 99.87\% against PPMI's own
three-class derivation, so the instability is a property of the definition, not
an error in the implementation.

\medskip
\noindent\textbf{Q. Why not just report one transition matrix?}\\
Because it is not a well-defined quantity. TD~\arr~PIGD is 0.100 at intervals
under 1.5 years and 0.289 beyond 4 years. A pooled matrix averages these in
whatever proportion PPMI's visit schedule and patient attendance happened to
produce, so it describes the study design as much as the disease.

\medskip
\noindent\textbf{Q. Which conversion definition should we use?}\\
The sustained one (two consecutive PIGD assessments) for anything presented as
a clinical outcome, because half of single-visit conversions do not hold. The
naive one is worth reporting alongside it to show the size of the effect. What
must not happen is choosing silently.

\medskip
\noindent\textbf{Q. Why Kaplan--Meier rather than counting conversions?}\\
Because follow-up length varies enormously --- median 3 visits, maximum 14
years. A raw proportion would conflate ``did not convert'' with ``was not
observed long enough''. KM censors correctly at the last labelled visit. The
curves are truncated at year 12 where the risk set drops below 25, since beyond
that each step is driven by a handful of patients.

\medskip
\noindent\textbf{Q. Why Wilson intervals?}\\
Row counts range from 334 to 2{,}261, and several probabilities sit near 0 or
1. The normal approximation produces intervals that fall outside $[0,1]$ in
exactly those cells; Wilson does not.

\medskip
\noindent\textbf{Q. Is PIGD an absorbing state, as the literature often assumes?}\\
Not in this data. 29.8\% of PIGD visits are followed by a non-PIGD assessment,
18.4\% of them straight back to TD. Some of that is measurement noise --- the
PIGD~\arr~TD reversion rate is 39.3\% --- but a model that forbids leaving PIGD
would be misspecified.

\medskip
\noindent\textbf{Q. You lose 23\% of visits. Does that break the analysis?}\\
It does not break it, but it biases it, and in a direction we can name. Roughly
half the loss is visits with no OFF-medication exam; the other half is exams
that were performed but could not be fully scored --- and in 1{,}243 of those
1{,}253 exams the missing item is postural stability, which is often unsafe to
test in a patient with poor balance. Those exams come from more advanced
patients (median Hoehn \& Yahr 3.0 versus 2.0, $p = 1.4\times10^{-17}$), and
postural stability is a PIGD item. So we preferentially lose the patients most
likely to be PIGD, which makes the reported drift toward PIGD a lower bound.
Section~\ref{sec:missingness} sets this out in full.

\medskip
\noindent\textbf{Q. Then why not relax the all-items rule?}\\
It was tested. Allowing one missing item per subscore recovers 44 visits and
changes zero existing labels, because the incomplete exams are typically
missing many items rather than one. The strict rule also matches PPMI's own
derivation, which is what lets us validate at 100\% against their
\code{td\_pigd}. The real remedy is imputation or an explicit
missing-not-at-random model, in Weeks 5--7.

\medskip
\noindent\textbf{Q. What did you actually deliver this fortnight?}\\
The four proposal tasks: subtype changes tracked per patient, transition
frequencies computed, probabilities estimated across four interval strata, and
pathways visualised. Plus two analyses that were not asked for --- reversion and
sustained conversion --- which changed how the first four should be read.
Outputs are 9 CSV tables, 6 figures and a text report, produced by one
reproducible command, with 74 unit tests over the underlying code.

\medskip
\noindent\textbf{Q. What is next?}\\
Weeks 5--7: the baseline probabilistic model. Given the interval dependence, a
continuous-time Markov chain estimating a generator $Q$ is a better baseline
than the discrete chain originally sketched, and it sets up the hierarchical
Bayesian stage directly.

% =============================================================================
\section{Reproducing this report}
% =============================================================================

\begin{verbatim}
    python scripts/00_data_audit.py          # verify source files
    python scripts/01_build_panel.py         # build labelled visit panel
    python scripts/02_transition_analysis.py # this report's numbers & figures
    python -m pytest tests/ -q               # 74 tests
\end{verbatim}

\noindent
Stage 2 writes to \code{outputs/transitions/}: nine CSV tables (transition
matrices with and without intervals, prevalence, reversion, pathways, both
Kaplan--Meier curves), six figures, and \code{report.txt} containing every
number quoted here. All item lists, thresholds, cohort rules and interval bins
live in \code{config/config.yaml}.

\vspace{1em}
\noindent\textbf{Reference.} Stebbins GT, Goetz CG, Burn DJ, Jankovic J, Khoo TK,
Tilley BC. How to identify tremor dominant and postural instability/gait
difficulty groups with the Movement Disorder Society Unified Parkinson's Disease
Rating Scale. \emph{Movement Disorders} 2013;28(5):668--70.

\end{document}
